\documentclass[10pt, aspectratio=169]{beamer}

\usepackage[utf8]{inputenc}
\usepackage[spanish, es-tabla]{babel}
\usepackage{booktabs}
\usepackage{array}
\usepackage{tabularx}
\usepackage{amsmath, amssymb}

\usetheme{Madrid}
\usecolortheme{beaver}

% Personalización institucional UCV
\setbeamercolor{title}{bg=darkred!80!black, fg=white}
\setbeamercolor{frametitle}{bg=darkred!80!black, fg=white}
\setbeamercolor{structure}{fg=darkred!80!black}

\title[Defensa de Anteproyecto --- EECA-UCV]{Relación Lineal entre el Tiempo de Uso Nocturno de Dispositivos Electrónicos y la Calidad del Sueño en Estudiantes de la EECA-UCV}
\subtitle{Anteproyecto de Investigación --- Período Lectivo 2026-01S}
\author[Armando, Giovanni, Anaís]{Armando Delfín \quad Giovanni Fermín \quad Anaís Moncada}
\institute[EECA - FACES - UCV]{
    Universidad Central de Venezuela\\
    Facultad de Ciencias Económicas y Sociales\\
    Escuela de Estadística y Ciencias Actuariales\\
    Cátedra de Computación Estadística\\
    \vspace{0.2cm}
    \textbf{Profesora:} Oriana Medina
}
\date[Caracas, Noviembre 2026]{Caracas, Noviembre de 2026}

\begin{document}

% -------------------------------------------------------------
% LÁMINA 1: PORTADA
% -------------------------------------------------------------
\begin{frame}
    \titlepage
\end{frame}

% -------------------------------------------------------------
% LÁMINA 2: PLANTEAMIENTO DEL PROBLEMA
% -------------------------------------------------------------
\begin{frame}{Planteamiento del Problema}
    \begin{itemize}
        \item \textbf{Fenómeno Global:} Adopción generalizada de dispositivos móviles con pantallas retroiluminadas LED y desplazamiento del horario de descanso en jóvenes universitarios.
        \item \textbf{Mecanismos Fisiológicos:}
        \begin{itemize}
            \item Emisión de luz azul ($\lambda \approx 450\text{--}480\text{ nm}$) que suprime la melatonina pineal y retrasa el ritmo circadiano maestro.
            \item Activación cognitiva e hiperalerta neuroconductual provocada por la interactividad digital.
            \item Desplazamiento deliberado de la hora de acostarse por interacción recreativa o académica.
        \end{itemize}
        \item \textbf{Contexto Local (EECA-UCV):} Alta exigencia analítica, horas prolongadas frente al computador y falta de evidencia cuantitativa propia sobre el déficit de descanso.
    \end{itemize}

    \vspace{0.3cm}
    \begin{alertblock}{Pregunta Central de Investigación}
        \centering
        \textbf{¿De qué manera la duración de la exposición nocturna a pantallas de dispositivos electrónicos altera la calidad y las horas efectivas de sueño en los estudiantes de la EECA-UCV?}
    \end{alertblock}
\end{frame}

% -------------------------------------------------------------
% LÁMINA 3: JUSTIFICACIÓN Y FINALIDAD
% -------------------------------------------------------------
\begin{frame}{Justificación y Finalidad}
    \begin{columns}[T]
        \begin{column}{0.48\textwidth}
            \begin{block}{Finalidad Institucional}
                \begin{itemize}
                    \item Dimensionar el problema latente del déficit de sueño en la comunidad universitaria.
                    \item Suministrar una \textbf{base métrica empírica objetiva} para que las autoridades de la EECA estructuren campañas de higiene del sueño fundamentadas en su propia población.
                    \item Sustituir las suposiciones genéricas sobre el cansancio estudiantil con datos directos.
                \end{itemize}
            \end{block}
        \end{column}
        \begin{column}{0.48\textwidth}
            \begin{block}{Aporte Metodológico}
                \begin{itemize}
                    \item Entrega de un \textbf{instrumento de medición validado} adaptado a la dinámica del estudiante venezolano.
                    \item Medición cronológica indirecta que elimina el sesgo de estimación subjetiva de horas en cama.
                    \item Precedente metodológico replicable en otras facultades para evaluar el ``jetlag social''.
                \end{itemize}
            \end{block}
        \end{column}
    \end{columns}
\end{frame}

% -------------------------------------------------------------
% LÁMINA 4: OBJETIVOS DE LA INVESTIGACIÓN
% -------------------------------------------------------------
\begin{frame}{Objetivos de la Investigación}
    \begin{block}{Objetivo General}
        \textbf{Analizar la relación lineal} entre el tiempo de uso nocturno de dispositivos electrónicos y la calidad del sueño en los estudiantes de la EECA-UCV durante el período 2026-01S.
    \end{block}

    \vspace{0.3cm}
    \begin{exampleblock}{Objetivos Específicos}
        \begin{enumerate}
            \item \textbf{Caracterizar} los hábitos de consumo digital nocturno según el tipo de dispositivo y la ventana de exposición previa al descanso.
            \item \textbf{Estimar} la duración absoluta y la calidad percibida del sueño durante los días lectivos frente a los fines de semana.
            \item \textbf{Identificar} las variaciones en los patrones de descanso asociadas a los distintos niveles de exposición a pantallas.
        \end{enumerate}
    \end{exampleblock}
\end{frame}

% -------------------------------------------------------------
% LÁMINA 5: DELIMITACIÓN Y COBERTURAS
% -------------------------------------------------------------
\begin{frame}{Delimitación: Coberturas y Períodos}
    \begin{columns}[T]
        \begin{column}{0.48\textwidth}
            \begin{block}{Cobertura Espacial y Poblacional}
                \begin{itemize}
                    \item \textbf{Cobertura Horizontal:} Circunscrita estrictamente al espacio físico y geográfico de la Facultad de Ciencias Económicas y Sociales (FACES-UCV), sede Caracas.
                    \item \textbf{Cobertura Vertical:} Desagregación analítica por sexo biológico, semestre cursado, tipo de dispositivo principal y tramos temporales de exposición nocturna.
                \end{itemize}
            \end{block}
        \end{column}
        \begin{column}{0.48\textwidth}
            \begin{block}{Delimitación Temporal}
                \begin{itemize}
                    \item \textbf{Período de Referencia:} Hábitos y comportamiento correspondientes al semestre académico activo 2026-01S.
                    \item \textbf{Período de Recolección de Datos:} Comprendido entre el 15 y el 30 de noviembre de 2026 (aplicación digital mediante Google Forms).
                \end{itemize}
            \end{block}
        \end{column}
    \end{columns}
\end{frame}

% -------------------------------------------------------------
% LÁMINA 6: SUSTENTO TEÓRICO Y BASES LEGALES
% -------------------------------------------------------------
\begin{frame}{Bases Teóricas y Respaldo Jurídico}
    \begin{columns}[T]
        \begin{column}{0.48\textwidth}
            \begin{block}{Modelo de Dos Procesos (Borbély, 1982)}
                \begin{itemize}
                    \item \textbf{Proceso C (Circadiano):} Marcapasos hipotalámico de 24h sensible a la luz azul.
                    \item \textbf{Proceso S (Homeostático):} Presión neuroquímica de sueño acumulada en vigilia.
                \end{itemize}
            \end{block}
        \end{column}
        \begin{column}{0.48\textwidth}
            \begin{block}{Marco Legal y Ético (CRBV, 1999)}
                \begin{itemize}
                    \item \textbf{Artículo 28 (Protección de Datos):} Derecho de acceso, reserva y fines informados (Consentimiento Informado).
                    \item \textbf{Artículo 60 (Intimidad y Vida Privada):} Blindaje jurídico del anonimato en hábitos personales.
                    \item \textbf{LFPE (2001) Arts. 19--20:} Secreto Estadístico estricto.
                \end{itemize}
            \end{block}
        \end{column}
    \end{columns}
\end{frame}

% -------------------------------------------------------------
% LÁMINA 7: OPERACIONALIZACIÓN DE VARIABLES
% -------------------------------------------------------------
\begin{frame}{Operacionalización de Variables}
    \centering
    \tiny
    \begin{tabular}{p{2.8cm}p{3.2cm}lll}
        \toprule
        \textbf{Variable} & \textbf{Definición Operacional} & \textbf{Rol Causal} & \textbf{Tipo} & \textbf{Escala} \\
        \midrule
        Tiempo de uso nocturno & Minutos en pantalla en las 2h previas al sueño & Independiente & Cuantitativa continua & Razón \\
        Dispositivo principal & Hardware predominante previo al descanso & Independiente & Cualitativa & Nominal \\
        Actividad principal & Tarea predominante en la pantalla & Independiente & Cualitativa & Nominal \\
        Latencia pre-sueño & Minutos entre apagar el equipo y dormir & Interviniente & Cuantitativa continua & Razón \\
        Horarios de sueño & Horas de acostarse y levantarse (Semana / Finde) & Primarias de cálculo & Cuantitativa & Intervalo \\
        Duración del sueño & Horas calculadas por diferencia horaria & Dependiente & Cuantitativa continua & Razón \\
        Interrupciones nocturnas & Conteo de despertares por el dispositivo & Interviniente & Cuantitativa discreta & Razón \\
        Calidad percibida & Valoración subjetiva al despertar por las mañanas & Dependiente & Cualitativa & Ordinal \\
        Variables de control & Sexo, Edad, Carrera EECA, Semestre & De control & Varias & Nom./Ord./Razón \\
        \bottomrule
    \end{tabular}
\end{frame}

% -------------------------------------------------------------
% LÁMINA 8: PLAN DE TABULACIONES BÁSICAS Y METODOLOGÍA
% -------------------------------------------------------------
\begin{frame}{Plan de Tabulaciones Básicas y Análisis Proyectado}
    \begin{columns}[T]
        \begin{column}{0.50\textwidth}
            \begin{block}{Esquema de Cuadros Proyectados (Vacíos)}
                \begin{itemize}
                    \item \textbf{Cuadro 1:} Distribución bivariada de dispositivos principales vs. tramos de tiempo de pantalla nocturno.
                    \item \textbf{Cuadro 2:} Estadísticos descriptivos ($\bar{x}, \tilde{x}, s, \text{IQR}$) de duración de sueño en días lectivos vs. fines de semana ($\Delta t_{\text{jetlag}}$).
                    \item \textbf{Cuadro 3:} Tabla de contingencia: Nivel de tiempo de pantalla vs. Calidad percibida del sueño.
                    \item \textbf{Cuadro 4:} Distribución de horas de sueño según sexo biológico.
                \end{itemize}
            \end{block}
        \end{column}
        \begin{column}{0.46\textwidth}
            \begin{alertblock}{Rigor Metodológico del Anteproyecto}
                \begin{itemize}
                    \item \textbf{Cero resultados preliminares:} Los cuadros se presentan completamente vacíos como matrices proyectadas.
                    \item \textbf{Asociación, no Causalidad:} Se contrasta la relación lineal ($r$ de Pearson) y tablas de contingencia ($\chi^2$), respetando la naturaleza observacional del estudio.
                    \item \textbf{Factibilidad Temporal:} Recolección y procesamiento diseñados para ejecutarse en 6 semanas.
                \end{itemize}
            \end{alertblock}
        \end{column}
    \end{columns}
\end{frame}

\end{document}

